% GENERATED FILE. Edit canonical lessons, then run npm run generate:books.
% canonical-source-hash: fnv1a64:3d240e6622b8a456
% canonical-lessons: ML-C32-undu, ML-C32-pokuka, ML-C32-varuka, ML-C32-tinnuka, ML-C32-kaanuka, ML-C32-ariyuka

\chapter{The Core Verbs}
\label{ch:ml-core-verbs}
\hlchaptermodality{32}

\begin{chapteropening}
\textbf{By the end of this chapter:} \emph{I can build a Malayalam verb as stem plus one ending with no person marking at all, swap that ending for past, future or mood, and say why 'I know Malayalam' puts me in the dative instead of the subject slot.}

Say \ml{എനിക്ക്} \ml{മലയാളം} \ml{അറിയാം}, name the two slots that built irikkunnu, pōkunnu and kāṇunnu, and explain why this verb moves the knower out of the subject slot altogether.
\end{chapteropening}

\section[uṇṭŭ]{\ml{ഉണ്ട്} (uṇṭŭ) — \textquotedblleft{}there is,\textquotedblright{} and the machine inside every Malayalam verb}

\label{lesson:ML-C32-undu}



\begin{quote}
Thirty-one chapters of words. Now the engine room. English makes one verb, \textquotedblleft{}to be,\textquotedblright{} do three jobs. Malayalam uses three words.
\end{quote}

\subsection*{What to know first}
\begin{quote}
\textbf{\ml{ഉണ്ട്}} — \emph{uṇṭŭ} — \textbf{there is, there are; {[}someone{]} has}
\end{quote}

\begin{quote}
\textbf{\ml{ഇരിക്കുക}} — \emph{irikkuka} — \textbf{to be somewhere, to stay, to sit}
\end{quote}

Chapter 2 gave you \textbf{\ml{ആണ്}} (\emph{āṇŭ}), the \textquotedblleft{}is\textquotedblright{} that links two nouns: \emph{enṟe pēru Arun āṇŭ}, \textquotedblleft{}my name is Arun.\textquotedblright{} \textbf{\ml{ഉണ്ട്}} does a different job — \textbf{existing} and \textbf{having}: \emph{enikku oru sahōdaran uṇṭŭ}, \textquotedblleft{}to-me one brother there-is,\textquotedblright{} which is how Malayalam says \textquotedblleft{}I have a brother.\textquotedblright{} \textbf{\ml{ഇരിക്കുക}} is being in a \textbf{place}: \emph{ñān vīṭṭil irikkunnu}, \textquotedblleft{}I am at home.\textquotedblright{}

Identity, existence, location — and none stands in for another.

\begin{sounds}
\textbf{\ml{ഉ}} (the standing vowel \emph{u}) + \textbf{\ml{ണ്}} + \textbf{\ml{ട്}} $\to$ \textbf{\ml{ഉണ്ട്}}. Both consonants are retroflex — tongue curled back — and both have had their vowel stripped, so the word trails off on a half-swallowed \emph{ŭ}: \emph{uṇ-ṭŭ}. \textbf{\ml{ഇരിക്കുക}} runs \emph{i-rik-ku-ka}, with the doubled \textbf{\ml{ക്ക}} held long.
\end{sounds}

\begin{cousinweb}
Both are native Dravidian, from the oldest layer, and the family split them between its daughters: \textbf{Tamil} kept \textbf{\ta{இரு}} (\emph{iru}), \textquotedblleft{}be, stay\textquotedblright{}; \textbf{Kannada} kept the same word, \textbf{\ka{ಇರು}} (\emph{iru}); \textbf{Telugu} kept the other one, \textbf{\te{ఉండు}} (\emph{uṇḍu}).

Malayalam kept \textbf{both}, and gave each a job.
\end{cousinweb}

\begin{grammarlens}[title={Grammar Lens: two pieces, where Tamil has three}]
English changes a verb with words \textbf{around} it: \emph{I am, I was, I will be}. Malayalam glues a piece \textbf{onto the end}:

\begin{quote}
\textbf{\ml{ഇരിക്ക്}} + \textbf{\ml{ഉന്നു}} $\to$ \textbf{\ml{ഇരിക്കുന്നു}}
\end{quote}

\begin{quote}
\emph{irikk-} (stay) + \emph{-unnu} (happening now) $\to$ \emph{irikkunnu}
\end{quote}

\textbf{Stem, then tense.} Two slots, and that is the whole machine.

Tamil glues on a \textbf{third} piece naming who acts: \emph{iru} + \emph{-kkiṟ-} + \emph{-ēṉ} (I) $\to$ \emph{irukkiṟēṉ}, \textquotedblleft{}I am,\textquotedblright{} against \emph{irukkiṟāṉ}, \textquotedblleft{}he is.\textquotedblright{} Malayalam threw that slot away. \emph{Ñān irikkunnu}, \emph{nī irikkunnu}, \emph{avan irikkunnu} — the verb never moves. Chapter 5 showed this with one verb; it holds for \textbf{every} Malayalam verb, and for no sister.
\end{grammarlens}

\subsection*{Your turn}
Say these aloud:

\begin{itemize}
\raggedright
  \item the three jobs — \textquotedblleft{}āṇŭ\textquotedblright{} is, \textquotedblleft{}uṇṭŭ\textquotedblright{} there-is, \textquotedblleft{}irikkuka\textquotedblright{} be somewhere
  \item \textquotedblleft{}enikku oru sahōdaran uṇṭŭ\textquotedblright{} — I have a brother
  \item the two pieces — \textquotedblleft{}irikk … unnu\textquotedblright{}
  \item one verb, three people — \textquotedblleft{}ñān, nī, avan … irikkunnu\textquotedblright{}
\end{itemize}

\begin{tcolorbox}[breakable,colback=teal!4,colframe=teal!35!black,title={Before you move on}]
Which be-word does Malayalam use for having something? (\textbf{\ml{ഉണ്ട്}}, \emph{uṇṭŭ}.) What are the two pieces of \emph{irikkunnu}? (\textbf{Stem and tense}; there is no third piece.) How does the verb change between \textquotedblleft{}I stay\textquotedblright{} and \textquotedblleft{}he stays\textquotedblright{}? (\textbf{It does not.}) Next: the piece that carries time.
\end{tcolorbox}
\section[pōkuka]{\ml{പോകുക} (pōkuka) — \textquotedblleft{}to go,\textquotedblright{} and the piece that carries time}

\label{lesson:ML-C32-pokuka}



\begin{quote}
Stem, then tense. Hold the stem still and swap the second piece. That is the whole of Malayalam tense, and it is a shorter list than you expect.
\end{quote}

\subsection*{What to know first}
\begin{quote}
\textbf{\ml{പോകുക}} — \emph{pōkuka} — \textbf{to go}
\end{quote}

Chapter 4 met this verb in passing, on the way to a farewell. Here it goes to work. The dictionary form is the stem \textbf{\ml{പോക്}-} (\emph{pōk-}) plus \textbf{-\ml{ഉക}} (\emph{-uka}), the ending that means \textquotedblleft{}to \_\_\_.\textquotedblright{}

\noindent
\begin{tabularx}{\linewidth}{@{}>{\raggedright\arraybackslash}X>{\raggedright\arraybackslash}X>{\raggedright\arraybackslash}X@{}}
\toprule
\textbf{when} & \textbf{Malayalam} & \textbf{the ending} \\
\midrule
now & \emph{pōkunnu} & \emph{-unnu} \\
already & \emph{pōyi} & \emph{-i} \\
not yet & \emph{pōkuṁ} & \emph{-uṁ} \\
\bottomrule
\end{tabularx}

You have said the middle one already: the goodbye \textbf{\ml{പോയി} \ml{വരാം}} (\emph{pōyi varāṁ}) opens with \emph{pōyi}, \textquotedblleft{}having gone.\textquotedblright{}

\begin{sounds}
\textbf{\ml{പോ}} (\emph{pa} wearing the long-\emph{ō} sign) + \textbf{\ml{കു}} (\emph{ku}) + \textbf{\ml{ക}} (\emph{ka}) $\to$ \textbf{\ml{പോകുക}}. The \emph{ō} is long and round; let it run.
\end{sounds}

\begin{cousinweb}
\textbf{\ml{പോ}-} is native Dravidian, so old it has barely moved: Tamil \textbf{\ta{போ}} (\emph{pō}), Telugu \textbf{\te{పోవు}} (\emph{pōvu}), the same verb.

Kannada's everyday word is \textbf{\ka{ಹೋಗು}} (\emph{hōgu}), and that \emph{h} is no different root. Old Dravidian \emph{p-} softened to \emph{h-} in Kannada — \textquotedblleft{}ten\textquotedblright{} is Malayalam \textbf{\ml{പത്ത്}} (\emph{pattŭ}) against Kannada \textbf{\ka{ಹತ್ತು}} (\emph{hattu}), the very same swap.
\end{cousinweb}

\begin{grammarlens}[title={Grammar Lens: three forms are the whole conjugation}]
In most languages a tense is a \textbf{list}. English has \emph{I go / he goes}. Tamil, Malayalam's closest sister, needs an ending for every person: \emph{pōgiṟēṉ} (\textquotedblleft{}I go\textquotedblright{}), \emph{pōgiṟāy} (\textquotedblleft{}you go\textquotedblright{}), \emph{pōgiṟāṉ} (\textquotedblleft{}he goes\textquotedblright{}), \emph{pōgiṟōm} (\textquotedblleft{}we go\textquotedblright{}).

Malayalam's present tense is one word long. \emph{Ñān pōkunnu}, \emph{nī pōkunnu}, \emph{avan pōkunnu}, \emph{avar pōkunnu} — I, you, he, they, and \emph{pōkunnu} every time. So those three forms are not the start of a conjugation. \textbf{They are the whole of it.}

That is the bargain. Nothing to memorise on the verb — and the pronoun in front can never be dropped, since it is now the only thing telling you who went.
\end{grammarlens}

\subsection*{Your turn}
Say these aloud:

\begin{itemize}
\raggedright
  \item \textquotedblleft{}pōkuka\textquotedblright{} — to go
  \item the three times — \textquotedblleft{}pōkunnu … pōyi … pōkuṁ\textquotedblright{}
  \item the pair that shares one sound-law — \textquotedblleft{}pattŭ … hattu\textquotedblright{}, \textquotedblleft{}pō … hōgu\textquotedblright{}
  \item four people, one verb — \textquotedblleft{}ñān, nī, avan, avar … pōkunnu\textquotedblright{}
\end{itemize}

\begin{tcolorbox}[breakable,colback=teal!4,colframe=teal!35!black,title={Before you move on}]
Say \textquotedblleft{}I will go.\textquotedblright{} (\emph{Ñān pōkuṁ}.) How many words are in this verb's whole Malayalam present tense? (\textbf{One} — \emph{pōkunnu}, for every person.) Is Kannada's \emph{hōgu} a different root from \emph{pō}? (\textbf{No} — the same old \emph{p} softened to \emph{h}.) Next: the one place a Malayalam verb can still surprise you.
\end{tcolorbox}
\section[varuka]{\ml{വരുക} (varuka) — \textquotedblleft{}to come,\textquotedblright{} and a past that reshapes the stem}

\label{lesson:ML-C32-varuka}



\begin{quote}
\emph{Pōkuka} took its endings without a fuss. Its lifelong partner does not — and the difference sits in exactly one of the three tenses.
\end{quote}

\subsection*{What to know first}
\begin{quote}
\textbf{\ml{വരുക}} — \emph{varuka} — \textbf{to come}
\end{quote}

The stem is \textbf{\ml{വരു}-} (\emph{varu-}), and the command is shorter still: \textbf{\ml{വാ}!} (\emph{vā!}), \textquotedblleft{}come!\textquotedblright{} (Chapter 4 met the dictionary form in its other common spelling, \textbf{\ml{വരിക}}, \emph{varika}; both name this verb.)

Go and come are inseparable here. The everyday goodbye \textbf{\ml{പോയി} \ml{വരാം}} (\emph{pōyi varāṁ}) is the two verbs holding hands — \textquotedblleft{}having gone, I'll come\textquotedblright{} — and Malayalam has no plain farewell that leaves the coming back out.

\begin{sounds}
\textbf{\ml{വ}} (\emph{va}) + \textbf{\ml{രു}} (\emph{ru}) + \textbf{\ml{ക}} (\emph{ka}) $\to$ \textbf{\ml{വരുക}}, three light beats. In the past form \textbf{\ml{വന്നു}} (\emph{vannu}) the \textbf{\ml{ന്ന}} is a doubled \emph{n}, held long: \emph{van-nu}, not \emph{vanu}.
\end{sounds}

\begin{cousinweb}
Native Dravidian, and every sister kept it — each reshaping the front of the word. \textbf{Tamil} has \textbf{\ta{வா}} (\emph{vā}) as the command and \textbf{\ta{வரு}-} (\emph{varu-}) as the stem; \textbf{Kannada} has \textbf{\ka{ಬಾ}} (\emph{bā}); \textbf{Telugu} has \textbf{\te{రా}} (\emph{rā}).

\emph{V}, \emph{b}, \emph{r}: three regular reworkings of one old word, not three verbs. And Malayalam's \emph{vā / varu-} pairing is Tamil's exactly, down to the two shapes — which is what \textquotedblleft{}closest sister\textquotedblright{} means in practice.
\end{cousinweb}

\begin{grammarlens}[title={Grammar Lens: where Malayalam keeps its irregularity}]
Line the two verbs up:

\noindent
\begin{tabularx}{\linewidth}{@{}>{\raggedright\arraybackslash}X>{\raggedright\arraybackslash}X>{\raggedright\arraybackslash}X@{}}
\toprule
\textbf{tense} & \textbf{go} & \textbf{come} \\
\midrule
now & \emph{pōkunnu} & \emph{varunnu} \\
already & \emph{pōyi} & \emph{vannu} \\
not yet & \emph{pōkuṁ} & \emph{varuṁ} \\
\bottomrule
\end{tabularx}

Present and future are built the same way for both: stem, plus \emph{-unnu} or \emph{-uṁ}. The past is not — \emph{varu-} collapses into \emph{vann-}, and no rule would have told you so.

That is Malayalam's whole irregularity budget, spent in one place. Once a verb's \textbf{past} is learned, nothing else about it can catch you out, because the missing person slot means there is no second, third or plural form waiting behind it. A Malayalam verb costs you three forms. A Tamil one costs three tenses times six persons.
\end{grammarlens}

\subsection*{Your turn}
Say these aloud:

\begin{itemize}
\raggedright
  \item the command — \textquotedblleft{}vā!\textquotedblright{}
  \item the three times — \textquotedblleft{}varunnu … vannu … varuṁ\textquotedblright{}
  \item the pair, past against past — \textquotedblleft{}pōyi … vannu\textquotedblright{}
  \item the goodbye you now understand — \textquotedblleft{}pōyi varāṁ\textquotedblright{}
\end{itemize}

\begin{tcolorbox}[breakable,colback=teal!4,colframe=teal!35!black,title={Before you move on}]
Which tense of \emph{varuka} is the irregular one? (\textbf{The past} — \emph{vannu}.) Say \textquotedblleft{}I will come.\textquotedblright{} (\emph{Ñān varuṁ}.) Are Kannada \emph{bā} and Telugu \emph{rā} different verbs from \emph{varuka}? (\textbf{No} — one old Dravidian word, reshaped three ways.) Next: a verb Malayalam kept where Tamil built a new one.
\end{tcolorbox}
\section[tinnuka]{\ml{തിന്നുക} (tinnuka) — \textquotedblleft{}to eat,\textquotedblright{} and the ending that names a verb's homeland}

\label{lesson:ML-C32-tinnuka}



\begin{quote}
Malayalam took in a great deal of Sanskrit. Eating is one of the places it never reached.
\end{quote}

\subsection*{What to know first}
\begin{quote}
\textbf{\ml{തിന്നുക}} — \emph{tinnuka} — \textbf{to eat}
\end{quote}

\begin{quote}
\textbf{\ml{ഉണ്ണുക}} — \emph{uṇṇuka} — \textbf{to eat a meal, to dine}
\end{quote}

Two inherited verbs, both alive. \emph{Tinnuka} is eating in general. \emph{Uṇṇuka} is sitting down to a proper meal, and gives the noun \textbf{\ml{ഊണ്}} (\emph{ūṇŭ}), \textquotedblleft{}lunch.\textquotedblright{} The tenses are regular: \emph{tinnunnu}, \emph{tinnu}, \emph{tinnuṁ}.

\begin{sounds}
\textbf{\ml{തി}} (\emph{ti}) + \textbf{\ml{ന്നു}} (a doubled \emph{nn} with \emph{u}) + \textbf{\ml{ക}} (\emph{ka}) $\to$ \textbf{\ml{തിന്നുക}}, with the long \emph{n} you held in \emph{vannu}. In \textbf{\ml{ഉണ്ണുക}} that doubled letter is retroflex \textbf{\ml{ണ്ണ}} (\emph{ṇṇ}): \emph{uṇ-ṇu-ka}.
\end{sounds}

\begin{cousinweb}
\textbf{\ml{തിന്നുക}} goes back to Dravidian \emph{tiṉ-}, and the family kept it: Telugu \textbf{\te{తిను}} (\emph{tinu}), Kannada \textbf{\ka{ತಿನ್ನು}} (\emph{tinnu}).

Tamil kept \emph{tiṉ} too, as \textbf{\ta{தின்}}, but narrowed it to chewing and built a \textbf{new} everyday word from two of its own: \textbf{\ta{சாப்பிடு}} (\emph{sāppiḍu}) is \emph{sāppu} (\textquotedblleft{}food\textquotedblright{}) plus \emph{iḍu} (\textquotedblleft{}to put\textquotedblright{}). Put food. Tamil's ordinary eat-verb is home-made; Malayalam's is the family's original.
\end{cousinweb}

\begin{grammarlens}[title={Grammar Lens: -\ml{ഉക} and -\ml{ഇക്കുക}, two doors into the verb}]
Every verb so far has ended in \textbf{-\ml{ഉക}} (\emph{-uka}): \emph{tinnuka}, \emph{pōkuka}, \emph{varuka}, \emph{irikkuka}. Chapter 5's \textquotedblleft{}to speak\textquotedblright{} did not — it was \textbf{\ml{സംസാരിക്കുക}} (\emph{saṁsārikkuka}), with a longer \textbf{-\ml{ഇക്കുക}} (\emph{-ikkuka}). That difference is a passport stamp:

\noindent
\begin{tabularx}{\linewidth}{@{}>{\raggedright\arraybackslash}X>{\raggedright\arraybackslash}X@{}}
\toprule
\textbf{ending} & \textbf{what it marks} \\
\midrule
\textbf{-\ml{ഉക}} \emph{-uka} & an inherited Dravidian verb \\
\textbf{-\ml{ഇക്കുക}} \emph{-ikkuka} & a Sanskrit word made into a verb \\
\bottomrule
\end{tabularx}

\emph{Saṁsārikkuka} is built on Sanskrit \emph{saṁsāra}; Chapter 9's \textbf{\ml{ക്ഷമിക്കുക}} (\emph{kṣamikkuka}, \textquotedblleft{}to forgive\textquotedblright{}) on \emph{kṣamā}. Malayalam needed one dependable way to make a verb of any borrowed noun, and \textbf{-\ml{ഇക്കുക}} is it: a naturalisation ceremony you can hear.
\end{grammarlens}

\subsection*{Your turn}
Say these aloud:

\begin{itemize}
\raggedright
  \item the two eat-verbs — \textquotedblleft{}tinnuka\textquotedblright{} (eat) and \textquotedblleft{}uṇṇuka\textquotedblright{} (dine)
  \item the three times — \textquotedblleft{}tinnunnu … tinnu … tinnuṁ\textquotedblright{}
  \item the cousins that kept the root — \textquotedblleft{}tinu … tinnu … tinnuka\textquotedblright{}
  \item the two endings — \textquotedblleft{}-uka is ours, -ikkuka came from Sanskrit\textquotedblright{}
\end{itemize}

\begin{tcolorbox}[breakable,colback=teal!4,colframe=teal!35!black,title={Before you move on}]
Did Malayalam borrow its everyday word for eating? (\textbf{No} — Tamil is the one that built a new compound.) Which verb is for sitting down to a meal? (\textbf{\ml{ഉണ്ണുക}}, \emph{uṇṇuka}.) What does \emph{-ikkuka} tell you? (\textbf{That the verb was made from a Sanskrit word.}) Next: a verb Malayalam kept where all three sisters walked away.
\end{tcolorbox}
\section[kāṇuka]{\ml{കാണുക} (kāṇuka) — \textquotedblleft{}to see,\textquotedblright{} and everything else the second slot can hold}

\label{lesson:ML-C32-kaanuka}



\begin{quote}
Two slots, and the second one has only carried tense so far. It holds a good deal more — and you have used one of its other endings since Chapter 4.
\end{quote}

\subsection*{What to know first}
\begin{quote}
\textbf{\ml{കാണുക}} — \emph{kāṇuka} — \textbf{to see, to meet}
\end{quote}

Present \emph{kāṇunnu}, future \emph{kāṇuṁ}, and a past that reshapes the way \emph{vannu} did: \textbf{\ml{കണ്ടു}} (\emph{kaṇṭu}), \textquotedblleft{}saw.\textquotedblright{}

You have said this verb in a goodbye. \textbf{\ml{നാളെ} \ml{കാണാം}} (\emph{nāḷe kāṇāṁ}) — \textquotedblleft{}see you tomorrow\textquotedblright{} — is \emph{nāḷe} (\textquotedblleft{}tomorrow\textquotedblright{}) plus this stem with a different ending.

\begin{sounds}
\textbf{\ml{കാ}} (\emph{ka} with the long-\emph{ā} sign) + \textbf{\ml{ണു}} (retroflex \emph{ṇu}) + \textbf{\ml{ക}} (\emph{ka}) $\to$ \textbf{\ml{കാണുക}}. That \textbf{\ml{ണ}} is the letter that ended \emph{uṇṭŭ}: tongue curled back to the roof of the mouth.
\end{sounds}

\begin{cousinweb}
\emph{Kāṇ-} is the family's own see-root, inherited from Proto-Dravidian — and here the four sisters part company. \textbf{Tamil} says \textbf{\ta{பார்}} (\emph{pār}), keeping \textbf{\ta{காண்}} (\emph{kāṇ}) for formal writing; \textbf{Kannada} says \textbf{\ka{ನೋಡು}} (\emph{nōḍu}), keeping \textbf{\ka{ಕಾಣು}} (\emph{kāṇu}) alongside; \textbf{Telugu} says \textbf{\te{చూడు}} (\emph{cūḍu}).

Malayalam simply kept the inherited word in daily use. Twice now — \emph{tinnuka} and \emph{kāṇuka} — the youngest of the four sisters turns out to be the conservative one.
\end{cousinweb}

\begin{grammarlens}[title={Grammar Lens: tense is only one thing that ending can be}]
Hold \emph{kāṇ-} still and change what follows it:

\noindent
\begin{tabularx}{\linewidth}{@{}>{\raggedright\arraybackslash}X>{\raggedright\arraybackslash}X@{}}
\toprule
\textbf{ending on \emph{kāṇ-}} & \textbf{what it says} \\
\midrule
\emph{kāṇunnu} & sees, is seeing \\
\emph{kāṇuṁ} & will see \\
\emph{kāṇāṁ} & may see; let's see \\
\emph{kāṇaṇaṁ} & must see \\
\emph{kāṇarutŭ} & must not see \\
\bottomrule
\end{tabularx}

Time, permission, obligation, prohibition — all in the same single slot. This is where the person ending went: Malayalam gave that slot up and spent the room on \textbf{mood} instead.

And the payoff lands on every line. \emph{Kāṇaṇaṁ} is \textquotedblleft{}I must see,\textquotedblright{} \textquotedblleft{}you must see\textquotedblright{} and \textquotedblleft{}they must see\textquotedblright{} without changing a letter. In Tamil each of those five ideas must still pick a person ending before it can be said aloud.
\end{grammarlens}

\subsection*{Your turn}
Say these aloud:

\begin{itemize}
\raggedright
  \item \textquotedblleft{}kāṇuka\textquotedblright{} — to see — and the past, \textquotedblleft{}kaṇṭu\textquotedblright{}
  \item the farewell you now understand — \textquotedblleft{}nāḷe kāṇāṁ\textquotedblright{}
  \item the ladder — \textquotedblleft{}kāṇunnu … kāṇuṁ … kāṇāṁ … kāṇaṇaṁ … kāṇarutŭ\textquotedblright{}
  \item one form, three people — \textquotedblleft{}ñān, nī, avar … kāṇaṇaṁ\textquotedblright{}
\end{itemize}

\begin{tcolorbox}[breakable,colback=teal!4,colframe=teal!35!black,title={Before you move on}]
Which of the four sisters kept the inherited see-word for everyday use? (\textbf{Malayalam} — the other three each chose a different word.) Say \textquotedblleft{}must see\textquotedblright{} and \textquotedblleft{}must not see.\textquotedblright{} (\emph{Kāṇaṇaṁ}; \emph{kāṇarutŭ}.) What besides tense does that second slot carry? (\textbf{Mood} — permission, obligation, prohibition.) Next: the verb that moves the person out of the sentence altogether.
\end{tcolorbox}
\section[aṟiyuka]{\ml{അറിയുക} (aṟiyuka) — \textquotedblleft{}to know,\textquotedblright{} the verb that takes no subject}

\label{lesson:ML-C32-ariyuka}



\begin{quote}
Five verbs, and every one let you stand in front of it as \emph{ñān}, \textquotedblleft{}I.\textquotedblright{} The last one will not — and Chapter 6 already handed you the sentence it makes.
\end{quote}

\subsection*{What to know first}
\begin{quote}
\textbf{\ml{അറിയുക}} — \emph{aṟiyuka} — \textbf{to know}
\end{quote}

Present \emph{aṟiyunnu}, past \textbf{\ml{അറിഞ്ഞു}} (\emph{aṟiññu}), and the form you have met: \textbf{\ml{അറിയാം}} (\emph{aṟiyāṁ}), with the \emph{-āṁ} ending of the last lesson.

Native Dravidian, and the sisters kept the root — Tamil \textbf{\ta{அறி}} (\emph{aṟi}), Kannada \textbf{\ka{ಅರಿ}} (\emph{ari}) — though for everyday knowing each reached elsewhere: Tamil \textbf{\ta{தெரியும்}} (\emph{teriyum}), Telugu \textbf{\te{తెలుసు}} (\emph{telusu}). Malayalam stayed with \emph{aṟi-}.

\begin{sounds}
\textbf{\ml{അ}} (the standing vowel \emph{a}) + \textbf{\ml{റി}} (\emph{ṟi}) + \textbf{\ml{യു}} (\emph{yu}) + \textbf{\ml{ക}} (\emph{ka}) $\to$ \textbf{\ml{അറിയുക}}. The \textbf{\ml{റ}} is Malayalam's hard, tapped \emph{ṟ}, a different letter and a different sound from the soft \textbf{\ml{ര}} (\emph{r}) of \emph{varuka}.
\end{sounds}

\begin{grammarlens}[title={Grammar Lens: the knower leaves the subject slot}]
\begin{quote}
\textbf{\ml{എനിക്ക്} \ml{മലയാളം} \ml{അറിയാം}.} — \emph{enikku malayāḷaṁ aṟiyāṁ} — \textquotedblleft{}I know Malayalam.\textquotedblright{}
\end{quote}

Chapter 6 gave you that sentence whole. Here is the verb inside it. You cannot say \emph{ñān malayāḷaṁ aṟiyāṁ}: the one who knows is not the subject. That person moves into the \textbf{dative}, the \textquotedblleft{}to me\textquotedblright{} case.

\begin{itemize}
\raggedright
  \item \emph{enikku aṟiyāṁ} — it is known \textbf{to me}
  \item \emph{ninakku aṟiyāṁ} — it is known \textbf{to you}
  \item \emph{avanŭ aṟiyāṁ} — it is known \textbf{to him}
\end{itemize}

Malayalam treats knowing as something that \textbf{arrives at} you rather than something you \textbf{do}, and marks it with the pronoun's case.

That is the chapter's closing symmetry. With the other five verbs the pronoun in front was the only thing telling you who acted, because the verb carried no person. Here it has changed case, so the sentence has no subject at all. Malayalam took the person off the verb; this construction takes it out of the subject slot too.
\end{grammarlens}

\subsection*{Your turn}
Say these aloud:

\begin{itemize}
\raggedright
  \item \textquotedblleft{}aṟiyuka\textquotedblright{} — to know — and the past, \textquotedblleft{}aṟiññu\textquotedblright{}
  \item \textquotedblleft{}enikku malayāḷaṁ aṟiyāṁ\textquotedblright{} — I know Malayalam
  \item three knowers, one verb — \textquotedblleft{}enikku … ninakku … avanŭ aṟiyāṁ\textquotedblright{}
  \item the chapter's six — \textquotedblleft{}uṇṭŭ, pōkunnu, varunnu, tinnunnu, kāṇunnu, aṟiyunnu\textquotedblright{}
\end{itemize}

\begin{tcolorbox}[breakable,colback=teal!4,colframe=teal!35!black,title={Before you move on}]
Say \textquotedblleft{}I know Malayalam,\textquotedblright{} and name the case the \textquotedblleft{}I\textquotedblright{} wears. (\emph{Enikku malayāḷaṁ aṟiyāṁ} — the \textbf{dative}.) Name the two slots of a Malayalam verb. (\textbf{Stem and tense}, with no third slot for person.) Which of the six verbs has an irregular past? (\emph{Varuka} $\to$ \emph{vannu}, and \emph{kāṇuka} $\to$ \emph{kaṇṭu}.) Which be-word says \textquotedblleft{}I have\textquotedblright{}? (\textbf{\ml{ഉണ്ട്}}, \emph{uṇṭŭ}.) How many forms has the Malayalam present tense, for any verb at all? (\textbf{One.})
\end{tcolorbox}
